\chapter{Prerequisites: what we assume and where to look}\label{ch:prereq}

This appendix lists, as precise statements without proofs, what the book assumes from the reader's
first courses: real analysis, linear algebra, complex analysis, the group theory, representation
theory and probability of Riley, Hobson and Bence \cite{RHB06}, and elementary statistics. Each
item is a card: a definition, or a theorem tagged \emph{(stated; reference)}, followed by the
chapter of the book where it is first needed and a reference where a proof can be found. Nothing
here is proved; the proofs are in the textbooks of \Cref{sec:prereq:books}. A reader who recognises
every card can skip the appendix. A reader who does not recognise one should read the corresponding
part of the reference before the chapter of the book that needs it.

The book assumes nothing beyond this appendix. In particular it does not assume abstract algebra
beyond the group theory of \Cref{sec:prereq:groups}: rings, fields, modules, Galois theory and
everything after them are built from the beginning in \Cref{ch:rings}, \Cref{ch:modules} and
\Cref{ch:galois}. Two warnings apply throughout. First, the reader's courses did linear algebra over
$\R$ and $\C$; every statement of \Cref{sec:prereq:linear} holds over an arbitrary field $k$ unless
the field is named, and the book uses them over $\F_q$, $\Q_\ell$, $\Q$ and finite extensions of
these. Second, the word ``character'' has two meanings in the book, RHB's (the trace function of a
representation) and the number-theoretic one (a one-dimensional representation, such as a Dirichlet
character or the cyclotomic character $\chil$); \Cref{sec:prereq:groups} says which is which.

\section{Real analysis}\label{sec:prereq:real}

This section is used from \Cref{ch:modules} on, where the ring $\Zl$ is constructed by the recipe
that constructs $\R$ from $\Q$, and in earnest in \Cref{ch:ellcurves}, \Cref{ch:modforms} and
\Cref{ch:lfunctions}, where lattice sums, $q$-expansions and Dirichlet series must be shown to
converge.

\begin{definition}[the real numbers; convergence]\label{def:prereq:completeness}
$\R$ is an ordered field with the \emph{least upper bound property}: every nonempty subset of $\R$
that is bounded above has a least upper bound (supremum) in $\R$. A sequence $(x_n)$ of real or
complex numbers \emph{converges} to $x$ if for every $\eps>0$ there is $N$ with $\abs{x_n-x}<\eps$
for all $n\ge N$; it is \emph{Cauchy} if for every $\eps>0$ there is $N$ with $\abs{x_m-x_n}<\eps$
for all $m,n\ge N$.
\end{definition}

\begin{theorem}[consequences of completeness (stated; reference)]\label{thm:prereq:completeness}
\begin{enumerate}[label=(\alph*)]
\item Every bounded monotone sequence of real numbers converges.
\item (Bolzano-Weierstrass) Every bounded sequence in $\R$ or $\C$ has a convergent subsequence.
\item (Cauchy criterion) A sequence in $\R$ or $\C$ converges if and only if it is Cauchy. In the
language of metric spaces: $\R$, $\C$ and $\R^n$ with the Euclidean distance are complete.
\item $\Q$ is not complete: a Cauchy sequence of rationals need not have a rational limit. $\R$ is,
up to isomorphism, the completion of $\Q$: the set of Cauchy sequences of rationals modulo the
sequences converging to $0$, with the operations performed termwise.
\end{enumerate}
\end{theorem}
\noindent\emph{First needed:} \Cref{ch:modules} and \Cref{ch:local}, where $\Zl$ and $\Q_p$ are
built from $\Z$ and $\Q$ by (d) with the $p$-adic absolute value in place of the usual one.
\emph{Reference:} \cite[Chapters 1 and 3, Theorem 3.11]{Rudin76}.

\emph{Two Cauchy sequences.} The decimal truncations $x_1=1.4$, $x_2=1.41$, $x_3=1.414,\dots$ of
$\sqrt2$ form a Cauchy sequence of rationals: for $m,n\ge N$ both $x_m$ and $x_n$ lie in the
interval $[x_N,x_N+10^{-N}]$, so $\abs{x_m-x_n}\le10^{-N}$. Its limit $x$ in $\R$ satisfies
$x^2=\lim x_n^2=2$, so the sequence has no limit in $\Q$; this is (d) in its simplest form. The
same recipe with another absolute value gives a different completion. Let $\abs{x}_3=3^{-v_3(x)}$,
where $v_3(x)$ is the exponent of $3$ in the rational number $x$, and let
$s_n=1+3+\dots+3^n=(3^{n+1}-1)/2$. For $m>n$ the difference $s_m-s_n$ is divisible by $3^{n+1}$, so
$\abs{s_m-s_n}_3\le3^{-n-1}$ and $(s_n)$ is Cauchy for $\abs{\cdot}_3$. In this case the limit is
even rational: $s_n+\tfrac12=3^{n+1}/2$ has $\abs{s_n+\tfrac12}_3=3^{-n-1}\to0$, so
$1+3+9+\dots=-\tfrac12$ is a true statement $3$-adically. Completing $\Q$ for $\abs{\cdot}_3$ gives
$\Q_3$; \Cref{ch:local} does this for every prime.

\begin{definition}[series]\label{def:prereq:series}
A series $\sum_{n\ge1}a_n$ of complex numbers \emph{converges} to $s$ if its partial sums
$\sum_{n\le N}a_n$ converge to $s$; it \emph{converges absolutely} if $\sum\abs{a_n}$ converges.
\end{definition}

\begin{theorem}[series (stated; reference)]\label{thm:prereq:series}
\begin{enumerate}[label=(\alph*)]
\item An absolutely convergent series converges, and $\abs{\sum a_n}\le\sum\abs{a_n}$.
\item (Comparison) If $\abs{a_n}\le b_n$ for all large $n$ and $\sum b_n$ converges, then
$\sum a_n$ converges absolutely.
\item (Ratio and root tests) If $\limsup\abs{a_{n+1}/a_n}<1$, or if $\limsup\abs{a_n}^{1/n}<1$,
then $\sum a_n$ converges absolutely; if $\limsup\abs{a_n}^{1/n}>1$ it diverges.
\item (Rearrangement) Every rearrangement of an absolutely convergent series converges absolutely
to the same sum.
\item (Double series) If the sums $\sum\abs{a_{m,n}}$ over finite sets of pairs $(m,n)$ are bounded,
then $\sum a_{m,n}$ taken in any order converges absolutely to the same value, which equals both
iterated sums $\sum_m\sum_n a_{m,n}$ and $\sum_n\sum_m a_{m,n}$. The same holds for sums indexed
by any countable set, such as the nonzero elements of a lattice in $\C$.
\item (Cauchy product) If $\sum a_n$ and $\sum b_n$ converge absolutely, then
$\sum_n\bigl(\sum_{k=0}^n a_kb_{n-k}\bigr)$ converges absolutely to the product of the two sums.
\item $\sum_{n\ge1}n^{-\sigma}$ converges for real $\sigma>1$ and diverges for $\sigma\le1$; for a
lattice $\Lambda\subset\C$, $\sum_{0\ne\omega\in\Lambda}\abs{\omega}^{-\sigma}$ converges if and
only if $\sigma>2$.
\end{enumerate}
\end{theorem}
\noindent\emph{First needed:} \Cref{ch:ellcurves} (the Weierstrass $\wp$-function and the Eisenstein
series $G_{2k}(\Lambda)=\sum'\omega^{-2k}$ are lattice sums, absolutely convergent for $2k\ge3$ by
(g) and summed in any order by (e)); \Cref{ch:lfunctions} (Dirichlet series and Euler products).
\emph{Reference:} \cite[Chapter 3]{Rudin76}; \cite[Chapter 8]{Apostol74}.

\emph{Why the lattice sum converges exactly for $\sigma>2$.} The second half of (g) reduces to the
first by counting lattice points. Let $\Lambda=\Z\omega_1+\Z\omega_2$ with $\omega_1,\omega_2$
linearly independent over $\R$. For an integer $r\ge1$, the points $m\omega_1+n\omega_2$ with
$\max(\abs m,\abs n)=r$ form the boundary of a square of lattice points, and there are
$(2r+1)^2-(2r-1)^2=8r$ of them. The function $(x,y)\mapsto\abs{x\omega_1+y\omega_2}$ is continuous
and nowhere zero (by linear independence) on the compact set $\max(\abs x,\abs y)=1$, so by
\Cref{thm:prereq:compact}(b) it has a minimum $c>0$ and a maximum $C$ there; by homogeneity every
point of the $r$-th square satisfies $cr\le\abs\omega\le Cr$. The contribution of the $r$-th square
to $\sum'\abs\omega^{-\sigma}$ therefore lies between $8C^{-\sigma}r^{1-\sigma}$ and
$8c^{-\sigma}r^{1-\sigma}$, and the whole sum converges if and only if $\sum_rr^{1-\sigma}$ does,
that is, if and only if $\sigma-1>1$. The terms are positive, so by (e) the order of summation does
not matter. In particular $G_{2k}(\Lambda)$ converges absolutely for $2k\ge4$, while the sum for
$G_2$ does not converge absolutely.

\begin{definition}[uniform convergence]\label{def:prereq:uniform}
Let $S$ be a set and $f_n,f\colon S\to\C$. Then $f_n\to f$ \emph{uniformly on $S$} if
$\sup_{x\in S}\abs{f_n(x)-f(x)}\to0$. A series $\sum f_n$ converges uniformly if its partial sums do.
For functions on an open subset $U$ of $\C$ or $\R^n$, \emph{locally uniform} convergence means
uniform convergence on every compact subset of $U$.
\end{definition}

\begin{theorem}[uniform convergence (stated; reference)]\label{thm:prereq:uniform}
\begin{enumerate}[label=(\alph*)]
\item A uniform limit of continuous functions on a metric space is continuous.
\item (Weierstrass $M$-test) If $\abs{f_n(x)}\le M_n$ for all $x\in S$ and $\sum M_n<\infty$, then
$\sum f_n$ converges absolutely and uniformly on $S$.
\item If $f_n\to f$ uniformly on $[a,b]$ and each $f_n$ is Riemann integrable, then $f$ is Riemann
integrable and $\int_a^bf_n\to\int_a^bf$; a uniformly convergent series may be integrated term by
term.
\item If the $f_n$ are differentiable on $[a,b]$, $f_n(x_0)$ converges for one point $x_0$, and
$f_n'$ converges uniformly, then $f_n$ converges uniformly to a differentiable $f$ with
$f'=\lim f_n'$. For holomorphic functions the much stronger \Cref{thm:prereq:weierstrass} holds.
\end{enumerate}
\end{theorem}
\noindent\emph{First needed:} \Cref{ch:ellcurves}; \Cref{ch:modforms} ($q$-expansions and
Eisenstein series converge locally uniformly on the upper half plane, by (b) on each compact set);
\Cref{ch:lfunctions}. \emph{Reference:} \cite[Chapter 7, Theorems 7.10 and 7.12]{Rudin76}.

\emph{The $M$-test for a $q$-expansion.} Suppose $\abs{a_n}\le Cn^k$ for all $n\ge1$ and put
$f(z)=\sum_{n\ge1}a_ne^{2\pi inz}$ on the upper half plane. A compact set $K$ in the upper half plane
has $\operatorname{Im}z\ge y_0$ on $K$ for some $y_0>0$ (the continuous function $\operatorname{Im}z$
attains a positive minimum on $K$, \Cref{thm:prereq:compact}(b)). On $K$,
$\abs{a_ne^{2\pi inz}}=\abs{a_n}e^{-2\pi n\operatorname{Im}z}\le Cn^ke^{-2\pi ny_0}=M_n$, and
$\sum M_n$ converges by the ratio test, since $M_{n+1}/M_n=(1+1/n)^ke^{-2\pi y_0}\to e^{-2\pi y_0}<1$.
So by (b) the series converges absolutely and uniformly on $K$: it converges locally uniformly on the
upper half plane, and \Cref{thm:prereq:weierstrass} then makes $f$ holomorphic. This is the argument
used, without further comment, for every $q$-expansion in the book.

\begin{definition}[power series]\label{def:prereq:power}
A power series is $\sum_{n\ge0}c_n(z-a)^n$ with $c_n,a\in\C$. Its \emph{radius of convergence} is
$R=1/\limsup\abs{c_n}^{1/n}\in[0,\infty]$.
\end{definition}

\begin{theorem}[power series (stated; reference)]\label{thm:prereq:power}
\begin{enumerate}[label=(\alph*)]
\item The series converges absolutely for $\abs{z-a}<R$, diverges for $\abs{z-a}>R$, and converges
uniformly on every closed disc $\abs{z-a}\le r$ with $r<R$.
\item Its sum $f$ is complex differentiable on the open disc of radius $R$, with derivative the
term-by-term derivative $\sum nc_n(z-a)^{n-1}$, which has the same radius of convergence.
\item The coefficients are determined by $f$: $c_n=f^{(n)}(a)/n!$. Two power series centred at
$a$ whose sums agree on a sequence of points $z_k\ne a$ converging to $a$ have the same
coefficients.
\item $\sum z^n/n!$ has $R=\infty$; $\sum z^n$ has $R=1$ and sum $1/(1-z)$.
\end{enumerate}
\end{theorem}
\noindent\emph{First needed:} \Cref{ch:modforms} (a $q$-expansion $\sum a_nq^n$ is a power series
in $q=e^{2\pi iz}$, $\abs{q}<1$). \emph{Warning:} the $p$-adic expansions $\sum a_np^n$ of
\Cref{ch:local} converge for a different absolute value, in which $\abs{p^n}\to0$; the analogy with
power series is useful, the theorems above do not transfer. \emph{Reference:} \cite[Chapter 3,
Theorem 3.39, and Chapter 8]{Rudin76}.

\emph{Example.} If $\abs{c_n}\le Cn^k$, then $\abs{c_n}^{1/n}\le C^{1/n}n^{k/n}\to1$, so the radius of
convergence is at least $1$: coefficients of polynomial growth, such as those of every $q$-expansion
in the book, always give a power series in $q$ converging on $\abs q<1$. The bound $1$ can be
attained: $\sum_{n\ge1}nq^n$ has radius exactly $1$, and differentiating $\sum q^n=1/(1-q)$ term by
term by (b) and multiplying by $q$ gives its sum $q/(1-q)^2$.

\begin{definition}[compactness]\label{def:prereq:compact}
A subset $K$ of $\R^n$ or $\C$ is \emph{compact} if every cover of $K$ by open sets has a finite
subcover; equivalently (Heine-Borel) if it is closed and bounded; equivalently if every sequence
in $K$ has a subsequence converging to a point of $K$.
\end{definition}

\begin{theorem}[continuity and compactness (stated; reference)]\label{thm:prereq:compact}
Let $K$ be compact and $f$ continuous on $K$.
\begin{enumerate}[label=(\alph*)]
\item $f(K)$ is compact.
\item If $f$ is real valued, it attains its maximum and its minimum on $K$.
\item $f$ is uniformly continuous on $K$.
\item A continuous bijection from a compact space onto a metric space is a homeomorphism.
\end{enumerate}
\end{theorem}
\noindent\emph{First needed:} \Cref{ch:curves} (a holomorphic function on a compact connected
Riemann surface is constant: (b) applied to $\abs{f}$ with the maximum modulus principle);
\Cref{ch:modforms} (fundamental domains, the compactification of $Y_0(N)$ by the cusps).
\emph{Reference:} \cite[Theorems 2.41, 4.14, 4.16 and 4.19]{Rudin76}.

\emph{How (b) is used on a Riemann surface.} Let $X$ be a compact connected Riemann surface and $f$
holomorphic on $X$. The continuous function $\abs f$ attains its maximum at some point $P$ by (b),
which holds for continuous real functions on any compact space. In a coordinate chart around $P$,
$f$ is a holomorphic function of one variable whose modulus has a local maximum, so it is constant
near $P$ by the maximum modulus principle (\Cref{thm:prereq:cauchy}(e)). The set of points of $X$
near which $f$ is identically $f(P)$ is open by definition, nonempty, and closed by the identity
theorem (\Cref{thm:prereq:cauchy}(d), applied in charts); since $X$ is connected it is all of $X$.
\Cref{ch:curves} gives the details.

\begin{definition}[the Riemann integral]\label{def:prereq:riemann}
For a bounded $f\colon[a,b]\to\R$ and a partition $a=x_0<\dots<x_n=b$, the upper and lower sums are
$\sum(x_i-x_{i-1})\sup_{[x_{i-1},x_i]}f$ and $\sum(x_i-x_{i-1})\inf_{[x_{i-1},x_i]}f$. The function
is \emph{Riemann integrable} if the infimum of the upper sums equals the supremum of the lower sums,
and $\int_a^bf$ is the common value. Complex-valued functions are integrated by real and imaginary
parts.
\end{definition}

\begin{theorem}[the Riemann integral (stated; reference)]\label{thm:prereq:riemann}
\begin{enumerate}[label=(\alph*)]
\item Every continuous function on $[a,b]$ is integrable, and so is every monotone function and
every bounded function with finitely many discontinuities.
\item The integral is linear and monotone, $\abs{\int_a^bf}\le\int_a^b\abs{f}$, and
$\int_a^b=\int_a^c+\int_c^b$.
\item (Fundamental theorem of calculus) If $f$ is continuous on $[a,b]$, then $F(x)=\int_a^xf$ is
differentiable with $F'=f$; if $F$ is differentiable on $[a,b]$ and $F'$ is integrable, then
$\int_a^bF'=F(b)-F(a)$.
\item Integration by parts and the substitution rule hold for $C^1$ functions.
\item For a $C^1$ curve $\gamma\colon[a,b]\to\C$ and $f$ continuous on its image, the line integral
is $\int_\gamma f(z)\,dz=\int_a^bf(\gamma(t))\gamma'(t)\,dt$; it does not depend on the
parametrisation of $\gamma$ and changes sign when the orientation is reversed.
\end{enumerate}
\end{theorem}
\noindent\emph{First needed:} \Cref{ch:curves} (line integrals define Cauchy's theorem and the
periods of a differential); \Cref{ch:ellcurves} (the periods of an elliptic curve over $\C$ are
integrals of $dx/y$ over closed paths); \Cref{ch:modforms} (the Petersson inner product is a double
integral over a fundamental domain). \emph{Reference:} \cite[Chapter 6, Theorems 6.8, 6.20 and
6.21]{Rudin76}.

\emph{The basic line integral.} Parametrise the unit circle by $\gamma(t)=e^{it}$, $0\le t\le2\pi$,
so $\gamma'(t)=ie^{it}$. For an integer $n$, (e) gives
\[
\int_{\abs z=1}z^n\,dz=\int_0^{2\pi}e^{int}\,ie^{it}\,dt=i\int_0^{2\pi}e^{i(n+1)t}\,dt
=\begin{cases}2\pi i&n=-1,\\0&n\ne-1,\end{cases}
\]
since $e^{i(n+1)t}$ has an antiderivative $e^{i(n+1)t}/(i(n+1))$ of period $2\pi$ when $n\ne-1$.
This one computation is behind the coefficient formula for Laurent series and the residue theorem
of \Cref{sec:prereq:complex}: integrating a Laurent series term by term around a circle keeps only
the coefficient of $z^{-1}$.

\begin{definition}[improper integrals]\label{def:prereq:improper}
$\int_a^\infty f=\lim_{b\to\infty}\int_a^bf$ when the limit exists; for $f$ unbounded near $0$,
$\int_0^1f=\lim_{\eps\to0^+}\int_\eps^1f$; $\int_0^\infty$ is the sum of $\int_0^1$ and
$\int_1^\infty$. The integral \emph{converges absolutely} if the corresponding integral of $\abs f$
converges.
\end{definition}

\begin{theorem}[improper integrals (stated; reference)]\label{thm:prereq:improper}
\leavevmode
\begin{enumerate}[label=(\alph*)]
\item (Comparison) Suppose that $\abs f\le g$ on $[a,\infty)$ and that $\int_a^\infty g$ converges. Then
$\int_a^\infty f$ converges absolutely; and absolutely convergent improper integrals converge.
\item $\int_1^\infty t^{-\sigma}\,dt$ converges if and only if $\sigma>1$;
$\int_0^1t^{\sigma-1}\,dt$ converges if and only if $\sigma>0$; $\int_0^\infty e^{-ct}t^m\,dt$
converges for every $c>0$ and $m>-1$.
\item (Parameters) If $f(t,s)$ is continuous on $[a,\infty)\times K$ with $K$ compact, and
$\abs{f(t,s)}\le g(t)$ with $\int_a^\infty g<\infty$, then $s\mapsto\int_a^\infty f(t,s)\,dt$ is
continuous on $K$; if moreover $\partial f/\partial s$ exists, is continuous and is dominated in the
same way, the integral may be differentiated under the integral sign.
\end{enumerate}
\end{theorem}

\begin{definition}[the Gamma function]\label{def:prereq:gamma}
For real $s>0$, $\Gamma(s)=\int_0^\infty e^{-t}t^{s-1}\,dt$. The integral converges absolutely: near
$0$ by \Cref{thm:prereq:improper}(b), and near $\infty$ because $e^{-t}t^{s-1}\le t^{-2}$ for
large $t$.
\end{definition}

\begin{theorem}[the Gamma function, real form (stated; reference)]\label{thm:prereq:gammareal}
\begin{enumerate}[label=(\alph*)]
\item $\Gamma(s+1)=s\,\Gamma(s)$ for $s>0$ (integration by parts).
\item $\Gamma(1)=1$, hence $\Gamma(n+1)=n!$ for integers $n\ge0$; $\Gamma(1/2)=\sqrt\pi$.
\item For $n\ge1$ and $s>0$, $\int_0^\infty e^{-nt}t^{s-1}\,dt=n^{-s}\Gamma(s)$ (substitute
$t\mapsto t/n$).
\item $\Gamma$ is continuous and positive on $(0,\infty)$, and $\log\Gamma$ is convex there.
\end{enumerate}
\end{theorem}
\noindent\emph{First needed:} \Cref{ch:lfunctions} (the completed $L$-function
$\Lambda(E,s)=N^{s/2}(2\pi)^{-s}\Gamma(s)L(E,s)$ and the Mellin transform of
\Cref{rem:prereq:mellin}); the complex form is \Cref{thm:prereq:gamma}. \emph{Reference:}
\cite[Chapter 8]{Rudin76}; \cite[Chapter 18]{RHB06}.

\emph{The computations behind (a) and (b), and behind the definition.} The bound
$e^{-t}t^{s-1}\le t^{-2}$ for large $t$ in \Cref{def:prereq:gamma} holds because
$t^{s+1}e^{-t}\to0$: choose an integer $m>s+1$; then $e^t\ge t^m/m!$ gives
$t^{s+1}e^{-t}\le m!\,t^{s+1-m}\to0$. For (a), integrate by parts on $[\eps,b]$:
\[
\int_\eps^be^{-t}t^s\,dt=\bigl[-e^{-t}t^s\bigr]_\eps^b+s\int_\eps^be^{-t}t^{s-1}\,dt,
\]
and the boundary terms $e^{-\eps}\eps^s$ and $e^{-b}b^s$ tend to $0$ as $\eps\to0^+$ (because $s>0$)
and $b\to\infty$. For (b), $\Gamma(1)=\int_0^\infty e^{-t}\,dt=1$, and induction on (a) gives
$\Gamma(n+1)=n\,\Gamma(n)=\dots=n!$. The substitution $t=u^2$, $dt=2u\,du$, turns $\Gamma(1/2)$ into
$2\int_0^\infty e^{-u^2}\,du=\int_{-\infty}^\infty e^{-u^2}\,du$, which is $\sqrt\pi$ (the
Gaussian integral, computed by squaring it and passing to polar coordinates).

With the analytic background in place, the next section turns to linear algebra, which the book
uses far more heavily and over fields the reader's courses did not consider.

\section{Linear algebra}\label{sec:prereq:linear}

This section is used from \Cref{ch:groups} on. Throughout, $k$ is a field and all vector spaces are
finite-dimensional over $k$; the reader's courses did $k=\R$ and $k=\C$, and the book uses
$k=\F_q$, $\Ql$, $\Q$ and their finite extensions, and the same constructions over the rings $\Z$
and $\Zl$ (\Cref{ch:modules} says what survives over a ring).

\begin{definition}[vector spaces, bases, dimension]\label{def:prereq:vs}
A \emph{vector space} over $k$ is an abelian group $V$ with a scalar multiplication $k\times V\to V$
satisfying the usual axioms. A family $(v_i)$ is \emph{linearly independent} if
$\sum a_iv_i=0$ forces all $a_i=0$, \emph{spanning} if every vector is a linear combination of it,
and a \emph{basis} if both. A \emph{linear map} $T\colon V\to W$ satisfies $T(av+bw)=aTv+bTw$; its
kernel and image are subspaces; $\End(V)$ is the ring of linear maps $V\to V$ and $\GL(V)$ its
group of units.
\end{definition}

\begin{theorem}[bases and dimension (stated; reference)]\label{thm:prereq:dimension}
\begin{enumerate}[label=(\alph*)]
\item Every finitely generated vector space has a basis, and any two bases have the same number of
elements, the \emph{dimension} $\dim V$. Over a finite field $\F_q$, a space of dimension $n$ has
$q^n$ elements.
\item Every linearly independent family extends to a basis; every spanning family contains a basis.
\item For a subspace $U\subseteq V$: $\dim U\le\dim V$, with equality if and only if $U=V$, and
$\dim V=\dim U+\dim(V/U)$.
\item (Rank-nullity) For a linear map $T\colon V\to W$, $\dim V=\dim\ker T+\dim\im T$. In
particular an injective, or a surjective, linear map between spaces of the same dimension is an
isomorphism.
\item $\dim(U+U')+\dim(U\cap U')=\dim U+\dim U'$ for subspaces $U,U'$ of $V$.
\end{enumerate}
\end{theorem}
\noindent\emph{First needed:} \Cref{ch:groups} (the lines of $\F_q^2$: there are $q+1$ of them,
the points of $\PP^1(\F_q)$); \Cref{ch:reps}. \emph{Reference:} \cite[Chapters 2 and 3]{Axler15};
\cite[Chapter 11]{DummitFoote}.

\emph{Counting lines.} A line (one-dimensional subspace) of $\F_q^2$ is spanned by any of its
nonzero vectors, and it contains $q-1$ of them by (a). Two distinct lines meet only in $0$, so the
$q^2-1$ nonzero vectors of $\F_q^2$ fall into $(q^2-1)/(q-1)=q+1$ lines. For $q=3$ these are the four
lines spanned by $(1,0)$, $(0,1)$, $(1,1)$, $(1,2)$.

\begin{definition}[matrices of linear maps; change of basis]\label{def:prereq:matrices}
Given bases $e=(e_1,\dots,e_n)$ of $V$ and $f=(f_1,\dots,f_m)$ of $W$, the \emph{matrix} of
$T\colon V\to W$ is the $m\times n$ matrix $A$ whose $j$-th column holds the coordinates of $Te_j$
in the basis $f$. Composition of maps corresponds to the product of matrices. If $e'=eP$ is another
basis of $V$ (the columns of the invertible $P$ are the coordinates of the $e'_j$ in $e$) and
$f'=fQ$ another basis of $W$, the matrix of $T$ in the new bases is $Q^{-1}AP$; for an
endomorphism and a single change of basis it is $P^{-1}AP$. The \emph{trace} of a square matrix
is the sum of its diagonal entries.
\end{definition}

The convention in \Cref{def:prereq:matrices} matters: books that record the old basis in terms of
the new one write $PAP^{-1}$ instead, with $P$ replaced by its inverse. A check of the book's
convention: in $V=k^2$ with the standard basis $e$, let $T$ have matrix $A=\sm1002$, and take the
new basis $e'_1=e_1$, $e'_2=e_1+e_2$, so $P=\sm1101$. Then
\[
P^{-1}AP=\begin{pmatrix}1&-1\\0&1\end{pmatrix}\begin{pmatrix}1&1\\0&2\end{pmatrix}
=\begin{pmatrix}1&-1\\0&2\end{pmatrix},
\]
and indeed $Te'_1=e_1=e'_1$ and $Te'_2=e_1+2e_2=-e'_1+2e'_2$, whose coordinates $(1,0)$ and
$(-1,2)$ are the columns found.

\begin{theorem}[determinants (stated; reference)]\label{thm:prereq:det}
The determinant $\det\colon\Mat_n(k)\to k$ is the unique function of the columns that is
multilinear, alternating, and takes the value $1$ at the identity matrix.
\begin{enumerate}[label=(\alph*)]
\item $\det(AB)=\det A\det B$ and $\det A^{\mathsf T}=\det A$.
\item $A$ is invertible if and only if $\det A\ne0$; then $A^{-1}=(\det A)^{-1}\operatorname{adj}A$,
where $\operatorname{adj}A$ is the transpose of the matrix of cofactors. Over a commutative ring the
same formula shows that $A$ is invertible if and only if $\det A$ is a unit.
\item The determinant of a block upper triangular matrix is the product of the determinants of
the diagonal blocks.
\item $\det(P^{-1}AP)=\det A$ and $\tr(P^{-1}AP)=\tr A$ (as $\tr(AB)=\tr(BA)$), so $\det T$ and
$\tr T$ are defined for an endomorphism $T$ independently of the basis.
\item $\det$ and $\tr$ are compatible with extension of the field: the determinant of $A$ viewed
over a larger field is the same element of $k$.
\end{enumerate}
\end{theorem}
\noindent\emph{First needed:} \Cref{ch:groups} ($\GL_2(\F_q)$, $\SL_2(\F_q)$ and their orders,
which are computed there); \Cref{ch:ellcurves} (the determinant of the Galois action on $E[m]$ is
the cyclotomic character, via the Weil pairing). \emph{Reference:} \cite[Chapter 10]{Axler15};
\cite[Chapter XIII]{Lang}.

\emph{Invertibility over a ring.} The ring version of (b) is the one the book uses over $\Z$ and
$\Zl$. The matrix $\sm2001$ has determinant $2$; its inverse $\sm{1/2}001$ has entries in $\Q$ and
in $\Zl$ for every odd $\ell$, where $2$ is a unit, but not in $\Z$ or $\Z_2$. So the same matrix is
in $\GL_2(\Z_3)$ and not in $\GL_2(\Z)$ or $\GL_2(\Z_2)$. By contrast $\sm2111$ has determinant $1$,
and its inverse $\sm1{-1}{-1}2$ has integer entries.

\begin{definition}[eigenvalues; characteristic and minimal polynomials]\label{def:prereq:eigen}
Let $T\in\End(V)$, $n=\dim V$. A scalar $\lambda\in k$ is an \emph{eigenvalue} of $T$ if
$\ker(T-\lambda)\ne0$; the nonzero elements of this kernel are the \emph{eigenvectors} and the
kernel is the \emph{eigenspace}. The \emph{characteristic polynomial} is
$\chi_T(X)=\det(X\cdot1-A)\in k[X]$ for any matrix $A$ of $T$; it is monic of degree $n$ and does
not depend on the basis. The \emph{minimal polynomial} $\mu_T$ is the monic polynomial of least
degree with $\mu_T(T)=0$.
\end{definition}

\begin{theorem}[eigenvalues (stated; reference)]\label{thm:prereq:eigen}
\begin{enumerate}[label=(\alph*)]
\item The eigenvalues of $T$ in $k$ are exactly the roots of $\chi_T$ in $k$.
\item $\chi_T(X)=X^n-(\tr T)X^{n-1}+\dots+(-1)^n\det T$; for $n=2$,
$\chi_T(X)=X^2-(\tr T)X+\det T$.
\item Over a field containing all the roots of $\chi_T$ (for instance an algebraic closure of $k$),
$\tr T$ is the sum and $\det T$ the product of the roots, counted with multiplicity; the
eigenvalues of $T^m$, and of $p(T)$ for a polynomial $p$, are the $m$-th powers, respectively the
values $p(\lambda)$, of the eigenvalues $\lambda$, with multiplicity.
\item (Cayley-Hamilton) $\chi_T(T)=0$. Consequently $\mu_T$ divides $\chi_T$; every polynomial
$p$ with $p(T)=0$ is a multiple of $\mu_T$; and $\mu_T$ and $\chi_T$ have the same irreducible
factors in $k[X]$, hence the same roots in an algebraic closure.
\item Eigenvectors belonging to distinct eigenvalues are linearly independent.
\end{enumerate}
\end{theorem}
\noindent\emph{First needed:} \Cref{ch:groups} (elements of $\GL_2(\F_3)$ of trace $0$ and
determinant $-1$); \Cref{ch:galrep} (the characteristic polynomial $X^2-a_pX+p$ of Frobenius on
$\Tl E$). \emph{Reference:} \cite[Chapters 5 and 8]{Axler15}; \cite[Chapter XIV]{Lang};
\cite[Section 12.2]{DummitFoote}.

\emph{One matrix over several fields.} The paper's matrix $J=\sm01{-1}0$ has trace $0$ and
determinant $1$, so $\chi_J=X^2+1$ by (b), and $J^2=-1$ by Cayley-Hamilton. Whether $J$ has
eigenvalues depends on the field. Over $\R$ it has none; over $\C$ it has $\pm i$. Over $\F_p$ with
$p$ odd it has eigenvalues exactly when $-1$ is a square modulo $p$: over $\F_5$ they are $2$ and $3$
($2^2=4=-1$), while over $\F_3$ the squares are $0$ and $1$, so there are none, and one must pass to
$\F_9=\F_3(i)$ to find $\pm i$. Over $\F_2$, $\chi_J=X^2+1=(X+1)^2$ has the double root $1$, but
$J=\sm0110\ne1$, so by \Cref{thm:prereq:jordan}(c) $J$ is conjugate to $\sm1101$ and is not
diagonalisable. Remark \ref{rem:prereq:changes} returns to the case of $\F_3$.

\begin{theorem}[diagonalisation and generalised eigenspaces (stated; reference)]\label{thm:prereq:diag}
Let $T\in\End(V)$, $n=\dim V$.
\begin{enumerate}[label=(\alph*)]
\item $T$ is \emph{diagonalisable} over $k$ (that is, $V$ has a basis of eigenvectors) if and only
if $\mu_T$ is a product of distinct linear factors in $k[X]$, if and only if $V$ is the direct sum
of the eigenspaces. In particular $T$ is diagonalisable if $\chi_T$ has $n$ distinct roots in $k$.
\item Suppose $k$ contains all the roots of $\chi_T$. Then $V=\bigoplus_\lambda V_\lambda$ with
$V_\lambda=\ker(T-\lambda)^n$ the \emph{generalised eigenspace}, each $V_\lambda$ is $T$-stable,
$\dim V_\lambda$ is the multiplicity of $\lambda$ as a root of $\chi_T$, and $T$ is diagonalisable
if and only if every $V_\lambda$ is the eigenspace $\ker(T-\lambda)$.
\item (Commuting operators) A family of pairwise commuting diagonalisable endomorphisms of $V$ is
simultaneously diagonalisable: $V$ has a basis of common eigenvectors. More generally, if $k$ is
algebraically closed, a family of pairwise commuting endomorphisms has a common eigenvector, and $V$
is the direct sum of the common generalised eigenspaces, each stable under the whole family.
\item $T$ is called \emph{semisimple} if it is diagonalisable over an algebraic closure of $k$,
\emph{nilpotent} if $T^m=0$ for some $m$, and \emph{unipotent} if $T-1$ is nilpotent. An
endomorphism of finite order $m$, with $m$ not divisible by the characteristic of $k$, is
semisimple with eigenvalues $m$-th roots of unity, because $\mu_T$ divides $X^m-1$, which has
distinct roots.
\end{enumerate}
\end{theorem}
\noindent\emph{First needed:} \Cref{ch:reps} (cyclic groups of order prime to the characteristic act
semisimply; groups of order divisible by the characteristic need not); \Cref{ch:hecke} (the Hecke
operators commute; the $T_n$ with $(n,N)=1$ are simultaneously diagonalisable on $S_2(\Gamma_0(N))$,
the $U_p$ need not be, and the generalised eigenspaces of the Hecke algebra are the objects of the
paper's Theorem thm:twist); \Cref{ch:modules} (generalised eigenspaces over a discrete valuation
ring). \emph{Reference:} \cite[Chapters 5 and 8]{Axler15}; \cite[Chapter XIV]{Lang}.

\emph{The hypothesis in (d) cannot be dropped.} Over $\F_\ell$ the matrix $u=\sm1101$ satisfies
$u^m=\sm1m01$, so it has order exactly $\ell$, the characteristic. Its minimal polynomial is
$(X-1)^2$ (as $u\ne1$ and $(u-1)^2=0$), which is not a product of distinct linear factors, so $u$ is
not diagonalisable over any field of characteristic $\ell$, although it has finite order. The
argument of (d) breaks exactly where expected: $X^\ell-1=(X-1)^\ell$ in characteristic $\ell$ has a
repeated root. Over $\Q$ the same matrix has infinite order. This $u$ is the shape of the inertia
action at a prime of type $\kod{IV}$ on $E[3]$ (\Cref{ch:galrep}).

\begin{theorem}[Jordan normal form (stated; reference)]\label{thm:prereq:jordan}
Suppose $k$ contains all the roots of $\chi_T$. Then $V$ has a basis in which the matrix of $T$ is
block diagonal with \emph{Jordan blocks} $J_m(\lambda)$: the $m\times m$ matrix with $\lambda$ on
the diagonal, $1$ just above the diagonal and $0$ elsewhere. The multiset of blocks is uniquely
determined by $T$, and $T$ is diagonalisable if and only if every block has size $1$. Consequences
for $\dim V=2$:
\begin{enumerate}[label=(\alph*)]
\item A unipotent $T\ne1$ is conjugate to $\sm1101$; its fixed space $\ker(T-1)$ is a line,
$(T-1)^2=0$, and over $k=\F_\ell$ it has order $\ell$. Likewise $-T$ (``minus a unipotent'') is
conjugate to $\sm{-1}{-1}0{-1}$ and has no nonzero fixed vector when the characteristic is odd.
\item If $\chi_T=(X-\lambda)(X-\mu)$ with $\lambda\ne\mu$ in $k$, then $T$ is diagonalisable with
eigenvalues $\lambda,\mu$.
\item If $\chi_T=(X-\lambda)^2$, then either $T=\lambda$ or $T$ is conjugate to $\sm\lambda10\lambda$.
\end{enumerate}
\end{theorem}
\noindent\emph{First needed:} \Cref{ch:groups} (unipotent elements of $\GL_2(\F_q)$ and the Borel
subgroup); \Cref{ch:reps}; \Cref{ch:galrep} (the image of inertia at a prime of type $\kod{IV}$ is
a nontrivial unipotent of order $3$ on $E[3]$, with a line of fixed points: the paper's Lemma
lem:IV and Remark rem:types); \Cref{ch:hecke} ($U_p$ with the repeated eigenvalue $0$ is a Jordan
block, not a scalar). \emph{Reference:} \cite[Chapter 8]{Axler15}; \cite[Section 12.3]{DummitFoote}.

\emph{The fixed vectors in (a), computed.} For $T=\sm1101$ and $v=(v_1,v_2)$ one has
$(T-1)v=(v_2,0)$, so $\ker(T-1)$ is the line spanned by $e_1$. The matrix $-T$ fixes $v$ if and only
if $(T+1)v=0$; but $T+1=\sm2102$ has determinant $4$, which is nonzero when the characteristic is
odd, so then $-T$ fixes only $0$. In characteristic $2$, $-T=T$ and the line of fixed points
returns.

\begin{definition}[bilinear forms; Gram matrices]\label{def:prereq:bilinear}
A \emph{bilinear form} on $V\times W$ is a map $B\colon V\times W\to k$ linear in each variable.
Its \emph{Gram matrix} with respect to bases $e$ of $V$ and $f$ of $W$ is $G=(B(e_i,f_j))$. The form
is \emph{perfect} (or nondegenerate) if $v\mapsto B(v,\cdot)$ is an isomorphism $V\to W^*=\Hom(W,k)$
and $w\mapsto B(\cdot,w)$ is an isomorphism $W\to V^*$; when $\dim V=\dim W$ each condition is
equivalent to $\det G\ne0$. For $V=W$ the form is \emph{symmetric} if $B(v,w)=B(w,v)$ and
\emph{alternating} if $B(v,v)=0$ for all $v$, which implies $B(v,w)=-B(w,v)$. A \emph{sesquilinear
form} on a complex space is linear in the first variable and conjugate-linear in the second, and
\emph{Hermitian} if $B(w,v)=\overline{B(v,w)}$.
\end{definition}

\begin{theorem}[Gram matrices under change of basis (stated; reference)]\label{thm:prereq:gram}
\leavevmode
\begin{enumerate}[label=(\alph*)]
\item If the bases change as $e'=eP$ and $f'=fQ$, then the Gram matrix in the new bases is $G'=P^{\mathsf T}GQ$. For $V=W$
and a single change of basis, $G'=P^{\mathsf T}GP$ and $\det G'=(\det P)^2\det G$. So the
determinant of the Gram matrix of a bilinear form on $V$ is well defined up to multiplication by
the square of a unit: of $k^\times$ over a field, of the unit group over a ring. Its $\ell$-adic
valuation, for a form on a free $\Zl$-module, is therefore an invariant of the form.
\item An alternating form on a two-dimensional space with basis $e_1,e_2$ is
$B(v,w)=a\det(v\mid w)$ with $a=B(e_1,e_2)$; its Gram matrix is $\sm0a{-a}0$ and its Gram
determinant is $a^2$, a square. More generally the Gram determinant of an alternating form on a
space of even dimension $2m$ is the square of a polynomial in the entries (the Pfaffian), and an
alternating form on a space of odd dimension is degenerate.
\item A symmetric bilinear form over a field of characteristic not $2$ has an orthogonal basis
(a basis in which $G$ is diagonal).
\end{enumerate}
\end{theorem}
\noindent\emph{First needed:} \Cref{ch:modules} (perfect pairings over $\Zl$; the Gram determinant
of a saturated sublattice and the order of the cokernel of $L\to\Hom(L,\mathcal O)$);
\Cref{ch:curves} (the cup product on $H^1$ of a curve is perfect and alternating);
\Cref{ch:congruence} (the pairing $\hat\delta_!$ restricted to the rank-two lattice $M_{f,\lambda}$
has Gram determinant $(\eta^\Sigma_{f,\lambda})^2$ up to a unit, by (b)). \emph{Reference:}
\cite[Chapter XV]{Lang}; \cite{Halmos74}.

\emph{Where (a) and (b) come from.} Write $e'_i=\sum_kP_{ki}e_k$ and $f'_j=\sum_lQ_{lj}f_l$. Bilinearity
gives $B(e'_i,f'_j)=\sum_{k,l}P_{ki}B(e_k,f_l)Q_{lj}=(P^{\mathsf T}GQ)_{ij}$, which is (a), and
$\det G'=\det P\det G\det Q$ by \Cref{thm:prereq:det}(a). For (b), an alternating form has
$B(e_1,e_1)=B(e_2,e_2)=0$ and $B(e_2,e_1)=-B(e_1,e_2)=-a$, so for $v=v_1e_1+v_2e_2$ and
$w=w_1e_1+w_2e_2$,
\[
B(v,w)=v_1w_2\,a+v_2w_1\,(-a)=a\,(v_1w_2-v_2w_1)=a\det(v\mid w).
\]
For example, the alternating form on $\Z_3^2$ with $a=3$ has Gram determinant $9$, of $3$-adic
valuation $2$, and no change of basis over $\Z_3$ can change that valuation, because $\det P$ is then
a unit of $\Z_3$.

\begin{paperlink}\label{pl:prereq:gram}
The paper's Lemma lem:ddt417 is stated for a free $\mathcal O$-module $L$ with a perfect bilinear
pairing for which the algebra $\T$ is self-adjoint, and $L[\wp]$ free of rank $d$ with Gram
determinant $\Delta$; its proof uses that the cokernel of $L[\wp]\to\Hom(L[\wp],\mathcal O)$ has
order $\#\mathcal O/\Delta$. In the application (Lemma lem:hecke and Step 2 of the proof of Theorem
thm:quant) $d=2$, the pairing is alternating, and $\Delta=(\eta)^2$ up to a unit by
\Cref{thm:prereq:gram}(b).
\end{paperlink}

\begin{definition}[inner products; adjoints]\label{def:prereq:inner}
An \emph{inner product} on a real vector space is a symmetric bilinear form with
$\langle v,v\rangle>0$ for $v\ne0$; on a complex vector space it is a Hermitian sesquilinear form
with the same positivity. It defines the norm $\norm v=\langle v,v\rangle^{1/2}$, orthogonality,
and the \emph{orthogonal complement} $U^\perp=\{v:\langle v,u\rangle=0\text{ for all }u\in U\}$ of
a subspace $U$. More generally, for a perfect bilinear form $B$ on $V$ (over any field, with no
positivity), the \emph{$B$-adjoint} of $T\in\End(V)$ is the unique $T^\dagger$ with
$B(Tv,w)=B(v,T^\dagger w)$ for all $v,w$, and $T$ is \emph{$B$-self-adjoint} if $T^\dagger=T$.
\end{definition}

\begin{theorem}[inner product spaces (stated; reference)]\label{thm:prereq:inner}
Let $V$ be a finite-dimensional real or complex inner product space.
\begin{enumerate}[label=(\alph*)]
\item (Cauchy-Schwarz) $\abs{\langle v,w\rangle}\le\norm v\norm w$. $V$ has an orthonormal basis
(Gram-Schmidt).
\item For a subspace $U$: $V=U\oplus U^\perp$, $(U^\perp)^\perp=U$ and $\dim U^\perp=\dim V-\dim U$.
\item Every $T\in\End(V)$ has a unique adjoint $T^*$ with $\langle Tv,w\rangle=\langle v,T^*w\rangle$;
its matrix in an orthonormal basis is the conjugate transpose of that of $T$; $(ST)^*=T^*S^*$; and a
subspace $U$ is $T$-stable if and only if $U^\perp$ is $T^*$-stable.
\item (Spectral theorem) If $T=T^*$, then $V$ has an orthonormal basis of eigenvectors of $T$ and
all eigenvalues are real; over $\C$ the same conclusion, with complex eigenvalues, holds for normal
$T$ ($TT^*=T^*T$).
\item A family of pairwise commuting self-adjoint (or normal) operators has a common orthonormal
basis of eigenvectors.
\item For a perfect bilinear form $B$ and a $B$-self-adjoint $T$, the $B$-orthogonal complement of
a $T$-stable subspace is $T$-stable, and $T$ need not be diagonalisable.
\end{enumerate}
\end{theorem}
\noindent\emph{First needed:} \Cref{ch:modforms} and \Cref{ch:hecke} (the Petersson inner product;
the $T_n$ with $(n,N)=1$ are self-adjoint for it, hence $S_2(\Gamma_0(N))$ has a basis of eigenforms
by (e); $U_p$ is not self-adjoint; $f^\perp$ is Hecke-stable by (c) because $f$ is an eigenvector of
the adjoints too); \Cref{ch:modules} (adjoints with respect to a pairing); \Cref{ch:moddeg} (the
adjoint of $T_n$ for the cup product is $w_NT_nw_N$); \Cref{ch:congruence}. \emph{Reference:}
\cite[Chapters 6 and 7]{Axler15}; \cite{Halmos74}.

\begin{definition}[dual spaces]\label{def:prereq:dual}
The \emph{dual} of $V$ is $V^*=\Hom_k(V,k)$. A basis $e_1,\dots,e_n$ of $V$ has the \emph{dual
basis} $e_1^*,\dots,e_n^*$ of $V^*$ with $e_i^*(e_j)=1$ if $i=j$ and $0$ otherwise. The
\emph{transpose} of $T\colon V\to W$ is $T^*\colon W^*\to V^*$, $\varphi\mapsto\varphi\circ T$. The
\emph{annihilator} of a subspace $U\subseteq V$ is $U^0=\{\varphi\in V^*:\varphi(U)=0\}$.
\end{definition}

\begin{theorem}[duality (stated; reference)]\label{thm:prereq:dual}
\begin{enumerate}[label=(\alph*)]
\item $\dim V^*=\dim V$, and the map $V\to V^{**}$, $v\mapsto(\varphi\mapsto\varphi(v))$, is an
isomorphism.
\item $\dim U^0=\dim V-\dim U$, $(U^0)^0=U$ under (a), $(V/U)^*=U^0$ and $U^*=V^*/U^0$.
\item The matrix of $T^*$ in the dual bases is the transpose of the matrix of $T$;
$\ker T^*=(\im T)^0$ and $\im T^*=(\ker T)^0$.
\item A perfect bilinear form $B\colon V\times W\to k$ identifies $W$ with $V^*$ and $V$ with $W^*$;
under this identification the $B$-adjoint of \Cref{def:prereq:inner} is the transpose.
\end{enumerate}
\end{theorem}
\noindent\emph{First needed:} \Cref{ch:modules} ($\Hom(L,\mathcal O)$ for a lattice $L$; Pontryagin
duality is the analogue for torsion modules); \Cref{ch:reps} (the contragredient representation
$V^*$); \Cref{ch:curves} ($H^1$ and $H_1$ of a curve are dual). \emph{Reference:}
\cite[Chapter 3]{Axler15}; \cite{Halmos74}.

\begin{definition}[tensor products]\label{def:prereq:tensor}
The \emph{tensor product} $V\otimes_kW$ is a vector space with a bilinear map
$V\times W\to V\otimes W$, $(v,w)\mapsto v\otimes w$, that is universal: every bilinear map
$V\times W\to U$ factors as a unique linear map $V\otimes W\to U$ composed with it. It exists and is
unique up to unique isomorphism. The \emph{symmetric square} $\Sym V$ is the quotient of
$V\otimes V$ by the span of the $v\otimes w-w\otimes v$, and the \emph{exterior square}
$\bigwedge^2V$ the quotient by the span of the $v\otimes v$.
\end{definition}

\begin{theorem}[tensor products (stated; reference)]\label{thm:prereq:tensor}
\begin{enumerate}[label=(\alph*)]
\item If $(e_i)$ is a basis of $V$ and $(f_j)$ a basis of $W$, then $(e_i\otimes f_j)$ is a basis
of $V\otimes W$; so $\dim(V\otimes W)=\dim V\cdot\dim W$. Every element of $V\otimes W$ is a finite
sum of elementary tensors, not in general a single one.
\item $\Hom(V,W)\iso V^*\otimes W$ via $(\varphi\otimes w)(v)=\varphi(v)w$. In particular
$\End(V)\iso V^*\otimes V$, and under this isomorphism the trace $\End(V)\to k$ is the contraction
$\varphi\otimes v\mapsto\varphi(v)$.
\item (Extension of scalars) For a field extension $k'/k$, $V_{k'}=V\otimes_kk'$ is a $k'$-vector
space with the same basis; a linear map $T$ extends to $T_{k'}$ with the same matrix, the same
characteristic and minimal polynomials, and the same trace and determinant.
\item (Functoriality) Linear maps $S\colon V\to V'$ and $T\colon W\to W'$ induce
$S\otimes T\colon V\otimes W\to V'\otimes W'$ with $(S\otimes T)(v\otimes w)=Sv\otimes Tw$; for
endomorphisms, $\tr(S\otimes T)=\tr S\cdot\tr T$, and over a field containing the eigenvalues
$\lambda_i$ of $S$ and $\mu_j$ of $T$ the eigenvalues of $S\otimes T$ are the products
$\lambda_i\mu_j$, with multiplicity.
\item $\dim\Sym V=n(n+1)/2$ and $\dim\bigwedge^2V=n(n-1)/2$ for $n=\dim V$. For $n=2$: $\Sym V$ has
dimension $3$, $\bigwedge^2V$ has dimension $1$ and $T$ acts on it by $\det T$; if $T$ has
eigenvalues $\lambda,\mu$ then it has eigenvalues $\lambda^2,\lambda\mu,\mu^2$ on $\Sym V$, and
$\lambda/\mu$, $1$, $1$, $\mu/\lambda$ on $V^*\otimes V=\End(V)$ (conjugation action), hence
$\lambda/\mu$, $1$, $\mu/\lambda$ on the trace-zero part when $\lambda\mu\ne0$.
\item The same definition, with ``bilinear'' meaning $R$-bilinear, gives the tensor product of
modules over a commutative ring $R$; (a) holds for free modules, (b) for finitely generated free
modules, and (d) in general. The modules $T\otimes_{\Zl}\Qlmod$ of the paper are of this kind.
\end{enumerate}
\end{theorem}
\noindent\emph{First needed:} \Cref{ch:reps} (tensor products of representations, the twist
$\rho\otimes\chi$ by a character, the adjoint representation $\End(V)=V^*\otimes V$ and its
trace-zero part $\ad\rho$, the identity $\ad\rho=\Sym\rho\otimes\det\rho^{-1}$ for
two-dimensional $\rho$); \Cref{ch:modules} (tensor products of modules); \Cref{ch:lfunctions}
(the eigenvalues $\alpha/\beta$, $1$, $\beta/\alpha$ of Frobenius on $\ad V$ at a good prime are
(e)); \Cref{ch:blochkato}. \emph{Reference:} \cite[Chapter XVI]{Lang};
\cite[Sections 10.4 and 11.5]{DummitFoote}; \cite{Halmos74}.

\begin{theorem}[the trace form (stated; reference)]\label{thm:prereq:traceform}
On $\Mat_n(k)$ the \emph{trace form} $(A,B)\mapsto\tr(AB)$ is a symmetric bilinear form, invariant
under conjugation ($\tr(gAg^{-1}\cdot gBg^{-1})=\tr(AB)$), and perfect: its Gram matrix with
respect to the elementary matrices $E_{ij}$ is a permutation matrix. On the trace-zero matrices
$\mathfrak{sl}_2(k)=\{A\in\Mat_2(k):\tr A=0\}$, with the basis
\[
e=\sm0100,\qquad h=\sm100{-1},\qquad f=\sm0010,
\]
one has $\tr(ef)=1$, $\tr(hh)=2$, and $\tr(ee)=\tr(ff)=\tr(eh)=\tr(fh)=0$; the Gram matrix with
respect to $(e,h,f)$ is
\[
\begin{pmatrix}0&0&1\\0&2&0\\1&0&0\end{pmatrix},\qquad\text{of determinant }-2.
\]
Hence the trace form on $\mathfrak{sl}_2$ is perfect over every field of characteristic not $2$,
and perfect on the free $\Zl$-module $\mathfrak{sl}_2(\Zl)$ for every odd prime $\ell$ (the
determinant $-2$ is a unit of $\Zl$); over $\F_2$ it is degenerate, $h$ being in its radical.
For $A\in\mathfrak{sl}_2(k)$ one has $\tr(A^2)=-2\det A$, and in general $\tr(A^2)=(\tr A)^2-2\det A$
for $A\in\Mat_2(k)$.
\end{theorem}
\noindent\emph{First needed:} \Cref{ch:reps} (the trace form makes $\ad\rho$ isomorphic to its own
dual when $2$ is invertible); \Cref{ch:blochkato} (for $A=\ad T$ and $\ell$ odd,
$T^D=\Hom(A,\Zl(1))\iso A(1)$ via the trace form, which the paper uses in Section 4.2);
\Cref{ch:modules}. \emph{Reference:} \cite[Chapter XV]{Lang}; the computation above is a direct
check.

\section{Complex analysis}\label{sec:prereq:complex}

This section is used from \Cref{ch:curves} on (Riemann surfaces, orders of meromorphic functions),
in \Cref{ch:ellcurves} (the Weierstrass $\wp$-function), in \Cref{ch:modforms} (the upper half
plane, $q$-expansions, the $\eta$-function) and in \Cref{ch:lfunctions} (Euler products, the Gamma
function, the Mellin transform). Throughout, $U$ is an open subset of $\C$.

\begin{definition}[holomorphic and meromorphic functions]\label{def:prereq:holo}
$f\colon U\to\C$ is \emph{holomorphic} if it is complex differentiable at every point of $U$;
\emph{entire} if holomorphic on $\C$; \emph{meromorphic} on $U$ if holomorphic on $U$ minus a
discrete set of points at each of which it has a pole (\Cref{thm:prereq:laurent}). A $C^1$ function
$u+iv$ is holomorphic if and only if $u_x=v_y$ and $u_y=-v_x$ (Cauchy-Riemann).
\end{definition}

\begin{theorem}[Cauchy; Taylor (stated; reference)]\label{thm:prereq:cauchy}
\begin{enumerate}[label=(\alph*)]
\item (Cauchy's theorem) If $f$ is holomorphic on $U$ and $\gamma$ is a closed piecewise $C^1$
curve in $U$ that can be contracted to a point inside $U$ (always the case when $U$ is a disc or is
simply connected), then $\int_\gamma f(z)\,dz=0$.
\item (Cauchy's integral formula) If $f$ is holomorphic on an open set containing the closed disc
$\abs{z-a}\le r$, then for $\abs{z-a}<r$ and all $n\ge0$,
$f^{(n)}(z)=\frac{n!}{2\pi i}\int_{\abs{w-a}=r}\frac{f(w)}{(w-z)^{n+1}}\,dw$.
\item (Taylor) If $f$ is holomorphic on the disc $\abs{z-a}<R$, then
$f(z)=\sum_{n\ge0}\frac{f^{(n)}(a)}{n!}(z-a)^n$ on that disc. Holomorphic functions are infinitely
differentiable, and (Morera) a continuous $f$ with $\int_{\partial\Delta}f=0$ for every triangle
$\Delta\subset U$ is holomorphic.
\item (Identity theorem) If $U$ is connected and $f,g$ holomorphic on $U$ agree on a set with an
accumulation point in $U$, then $f=g$. The zeros of a nonconstant holomorphic function on a
connected $U$ are isolated, and each has a finite \emph{order} $\ord_af$, the least $n$ with
$f^{(n)}(a)\ne0$; then $f(z)=(z-a)^{\ord_af}g(z)$ with $g$ holomorphic and $g(a)\ne0$.
\item (Liouville) A bounded entire function is constant. (Maximum modulus) If $U$ is connected and
$\abs f$ has a local maximum in $U$, then $f$ is constant.
\end{enumerate}
\end{theorem}
\noindent\emph{First needed:} \Cref{ch:curves}. \emph{Reference:} \cite[Chapter 4]{Ahlfors79};
\cite[Chapters 2 and 3]{SteinShakarchi03}.

\begin{theorem}[Laurent series; poles and residues (stated; reference)]\label{thm:prereq:laurent}
If $f$ is holomorphic on the annulus $r<\abs{z-a}<R$, then
\[
f(z)=\sum_{n\in\Z}c_n(z-a)^n,\qquad c_n=\frac1{2\pi i}\int_{\abs{w-a}=s}f(w)(w-a)^{-n-1}\,dw
\quad(r<s<R),
\]
the series converging absolutely and locally uniformly on the annulus, and the $c_n$ are unique.
When $r=0$ (an \emph{isolated singularity} at $a$) exactly one of the following holds: $c_n=0$ for
all $n<0$ and $f$ extends holomorphically to $a$ (\emph{removable}; equivalently $f$ is bounded
near $a$); there is $m\ge1$ with $c_{-m}\ne0$ and $c_n=0$ for $n<-m$ (a \emph{pole of order $m$};
equivalently $\abs{f(z)}\to\infty$ as $z\to a$); or neither (\emph{essential}). In the first two
cases $\ord_af$ is the least $n$ with $c_n\ne0$ (negative at a pole), and
$\ord_a(fg)=\ord_af+\ord_ag$, $\ord_a(f+g)\ge\min(\ord_af,\ord_ag)$. The \emph{residue} is
$\operatorname{Res}_af=c_{-1}$; at a simple pole it is $\lim_{z\to a}(z-a)f(z)$. The meromorphic
functions on a connected $U$ form a field.
\end{theorem}
\noindent\emph{First needed:} \Cref{ch:curves} (the order of a meromorphic function at a point of a
Riemann surface, and the divisor of a function); \Cref{ch:modforms} (the $q$-expansion at a cusp is
a Laurent series in $q$; a cusp form has no terms of order $\le0$). \emph{Reference:}
\cite[Chapter 5, Section 1]{Ahlfors79}; \cite[Chapter 3]{SteinShakarchi03}.

\begin{theorem}[the residue theorem (stated; reference)]\label{thm:prereq:residue}
Let $f$ be holomorphic on $U$ except at isolated points $a_1,\dots,a_k$, and let $\gamma$ be a
closed piecewise $C^1$ curve in $U$, not passing through any $a_j$, contractible in $U$. Then
\[
\frac1{2\pi i}\int_\gamma f(z)\,dz=\sum_jn(\gamma,a_j)\operatorname{Res}_{a_j}f,
\]
where $n(\gamma,a)$ is the winding number of $\gamma$ about $a$ ($1$ for a simple closed curve
traversed counterclockwise with $a$ inside, $0$ with $a$ outside). (Argument principle) If $f$ is
meromorphic on $U$ and has no zeros or poles on a simple closed curve $\gamma$, then
$\frac1{2\pi i}\int_\gamma f'/f$ equals the number of zeros minus the number of poles inside
$\gamma$, counted with multiplicity, that is, $\sum_a\ord_af$ over the points inside.
\end{theorem}
\noindent\emph{First needed:} \Cref{ch:curves} (a meromorphic function on a compact Riemann surface
has as many zeros as poles: the degree of a principal divisor is $0$); \Cref{ch:modforms} (the
valence formula is the argument principle on a fundamental domain). \emph{Reference:}
\cite[Chapter 4, Section 5]{Ahlfors79}; \cite[Chapter 3]{SteinShakarchi03}.

\begin{theorem}[Weierstrass: uniform limits of holomorphic functions (stated; reference)]\label{thm:prereq:weierstrass}
If $f_n$ are holomorphic on $U$ and $f_n\to f$ uniformly on every compact subset of $U$, then $f$
is holomorphic on $U$ and $f_n^{(k)}\to f^{(k)}$ uniformly on every compact subset, for every
$k\ge1$. Consequently a series $\sum f_n$ of holomorphic functions that converges locally uniformly
(for instance by the $M$-test on each compact set) has a holomorphic sum and may be differentiated
term by term.
\end{theorem}
\noindent\emph{First needed:} \Cref{ch:ellcurves} ($\wp(z)=z^{-2}+\sum'\bigl((z-\omega)^{-2}-\omega^{-2}\bigr)$
is meromorphic with double poles at the lattice points); \Cref{ch:modforms} (Eisenstein series and
$q$-expansions define holomorphic functions on the upper half plane); \Cref{ch:lfunctions} (a
Dirichlet series is holomorphic in its half-plane of absolute convergence). \emph{Reference:}
\cite[Chapter 5, Section 1]{Ahlfors79}; \cite[Chapter 2]{SteinShakarchi03}.

\begin{definition}[the Riemann sphere; Mobius transformations]\label{def:prereq:mobius}
The \emph{Riemann sphere} is $\widehat\C=\C\cup\{\infty\}$, with the neighbourhoods of $\infty$
the complements of compact subsets of $\C$ and the coordinate $w=1/z$ near $\infty$; a function is
holomorphic or meromorphic at $\infty$ if $f(1/w)$ is so at $w=0$. A \emph{Mobius transformation}
is $z\mapsto(az+b)/(cz+d)$ with $a,b,c,d\in\C$ and $ad-bc\ne0$, extended to $\widehat\C$ by
$-d/c\mapsto\infty$ and $\infty\mapsto a/c$ (and $\infty\mapsto\infty$ if $c=0$).
\end{definition}

\begin{theorem}[Mobius transformations (stated; reference)]\label{thm:prereq:mobius}
\begin{enumerate}[label=(\alph*)]
\item Under composition the Mobius transformations form a group, and the map
$\sm abcd\mapsto\bigl(z\mapsto\frac{az+b}{cz+d}\bigr)$ is a surjective homomorphism from
$\GL_2(\C)$ onto it with kernel the scalar matrices; so the group is
$\PGL_2(\C)=\GL_2(\C)/\C^\times=\SL_2(\C)/\{\pm1\}$.
\item The action on $\widehat\C$ is sharply triply transitive: for any two triples of distinct
points there is exactly one Mobius transformation taking the first to the second.
\item Mobius transformations map the family of circles and lines in $\C$ (lines being the circles
through $\infty$) to itself.
\item A Mobius transformation $g\ne1$ with matrix in $\SL_2(\C)$ has two fixed points on
$\widehat\C$ if $(\tr g)^2\ne4$ and one if $(\tr g)^2=4$.
\item The meromorphic functions on $\widehat\C$ are exactly the rational functions, and the
biholomorphic maps $\widehat\C\to\widehat\C$ are exactly the Mobius transformations.
\end{enumerate}
\end{theorem}
\noindent\emph{First needed:} \Cref{ch:curves} ($\PP^1(\C)=\widehat\C$ is the first compact Riemann
surface; the degree of a rational function); \Cref{ch:modforms} (the action of $\SL_2(\Z)$ on the
cusps $\PP^1(\Q)=\Q\cup\{\infty\}$). \emph{Reference:} \cite[Chapter 3, Section 3]{Ahlfors79};
\cite[Chapter 8]{SteinShakarchi03}.

\begin{theorem}[the upper half plane (stated; reference)]\label{thm:prereq:halfplane}
Let $\HH=\{z\in\C:\operatorname{Im}z>0\}$ and, for $g=\sm abcd\in\SL_2(\R)$, $gz=(az+b)/(cz+d)$.
\begin{enumerate}[label=(\alph*)]
\item $\operatorname{Im}(gz)=\operatorname{Im}(z)/\abs{cz+d}^2$, so $g$ maps $\HH$ to itself,
$(gh)z=g(hz)$, and $\SL_2(\R)$ acts on $\HH$; the kernel of the action is $\{\pm1\}$. More
generally $\GL_2(\R)^+$, the matrices of positive determinant, acts by the same formula, with
$\operatorname{Im}(gz)=\det(g)\operatorname{Im}(z)/\abs{cz+d}^2$.
\item The action is transitive, and the stabiliser of $i$ is
$\mathrm{SO}_2(\R)=\{\sm{\cos\theta}{\sin\theta}{-\sin\theta}{\cos\theta}\}$.
\item Every biholomorphic map $\HH\to\HH$ is $z\mapsto gz$ for some $g\in\SL_2(\R)$, unique up to
sign: $\Aut(\HH)=\mathrm{PSL}_2(\R)$. The Cayley map $z\mapsto(z-i)/(z+i)$ is a biholomorphic map of
$\HH$ onto the unit disc.
\item $\frac{d}{dz}(gz)=(cz+d)^{-2}$, and the measure $dx\,dy/y^2$ on $\HH$ ($z=x+iy$) is invariant
under $\SL_2(\R)$.
\item An element $g\ne\pm1$ of $\SL_2(\R)$ is \emph{elliptic} if $\abs{\tr g}<2$ (one fixed point in
$\HH$), \emph{parabolic} if $\abs{\tr g}=2$ (one fixed point, in $\R\cup\{\infty\}$), and
\emph{hyperbolic} if $\abs{\tr g}>2$ (two fixed points in $\R\cup\{\infty\}$). The group also acts
on $\R\cup\{\infty\}$, and $\SL_2(\Z)$ acts on $\Q\cup\{\infty\}$.
\end{enumerate}
\end{theorem}
\noindent\emph{First needed:} \Cref{ch:ellcurves} (lattices in $\C$ up to homothety correspond to
the points of $\SL_2(\Z)\backslash\HH$); \Cref{ch:modforms} (everything: the modular group,
fundamental domains, cusps, elliptic points, the Petersson inner product via (d)).
\emph{Reference:} \cite[Chapter 3, Section 3]{Ahlfors79}; \cite[Chapter 1]{DiamondShurman}.

\begin{theorem}[the exponential and the logarithm (stated; reference)]\label{thm:prereq:exp}
\leavevmode
\begin{enumerate}[label=(\alph*)]
\item The function $\exp(z)=\sum_{n\ge0}z^n/n!$ is entire, $\exp(z+w)=\exp z\exp w$, $\exp'=\exp$,
$\abs{\exp(x+iy)}=e^x$, and $\exp\colon\C\to\C^\times$ is surjective with kernel $2\pi i\Z$.
\item The map $z\mapsto q=e^{2\pi iz}$ sends $\HH$ onto the punctured disc $0<\abs q<1$, with
$q(z)=q(z')$ if and only if $z-z'\in\Z$, and $\abs q=e^{-2\pi\operatorname{Im}z}\to0$ as
$\operatorname{Im}z\to\infty$. Consequently a holomorphic $f$ on $\HH$ with $f(z+1)=f(z)$ is
$f(z)=F(q)$ for a unique holomorphic $F$ on the punctured disc, and by \Cref{thm:prereq:laurent}
\[
f(z)=\sum_{n\in\Z}a_nq^n,\qquad a_n=\int_0^1f(x+iy)\,e^{-2\pi in(x+iy)}\,dx\quad\text{for any }y>0,
\]
converging absolutely and locally uniformly on $\HH$: this is the \emph{$q$-expansion}, or Fourier
expansion, of $f$, and ``holomorphic at $\infty$'' means $a_n=0$ for $n<0$.
\item On a simply connected open set not containing $0$ there is a holomorphic branch of the
logarithm, unique up to adding a multiple of $2\pi i$; the principal branch on
$\C\setminus(-\infty,0]$ is $\operatorname{Log}(re^{i\theta})=\log r+i\theta$ with $-\pi<\theta<\pi$,
and $z^s=\exp(s\operatorname{Log}z)$ for $s\in\C$. For a positive integer $n$ and $s\in\C$,
$n^{-s}=\exp(-s\log n)$ with the real logarithm, and $\abs{n^{-s}}=n^{-\operatorname{Re}s}$.
\item $q^{1/24}=e^{2\pi iz/24}$ is a well-defined holomorphic function on $\HH$ (a function of
$z$, not of $q$); it is the factor in front of the $\eta$-function.
\end{enumerate}
\end{theorem}
\noindent\emph{First needed:} \Cref{ch:modforms}; \Cref{ch:lfunctions} (Dirichlet series).
\emph{Reference:} \cite[Chapter 2, Section 3]{Ahlfors79}; \cite[Chapter 1]{DiamondShurman}.

\begin{theorem}[infinite products (stated; reference)]\label{thm:prereq:products}
\begin{enumerate}[label=(\alph*)]
\item If $\sum\abs{a_n}<\infty$, then the partial products of $\prod_{n\ge1}(1+a_n)$ converge, the
limit is $0$ if and only if some factor is $0$, and the limit does not depend on the order of the
factors. (One says the product converges absolutely.)
\item If the $a_n$ are holomorphic on $U$ and $\sum\abs{a_n}$ converges uniformly on every compact
subset of $U$, then $P=\prod(1+a_n)$ is holomorphic on $U$, $P(z_0)=0$ if and only if some
$1+a_n(z_0)=0$, $\ord_{z_0}P=\sum_n\ord_{z_0}(1+a_n)$, and where $P\ne0$ its logarithmic derivative
is $P'/P=\sum a_n'/(1+a_n)$.
\item For $\abs q<1$ the product $\prod_{n\ge1}(1-q^n)$ converges absolutely, uniformly on
$\abs q\le r<1$, to a nonvanishing holomorphic function of $q$; so
$\eta(z)=q^{1/24}\prod_{n\ge1}(1-q^n)$, $q=e^{2\pi iz}$, is holomorphic and nonvanishing on $\HH$.
\item For $\operatorname{Re}s>1$, $\prod_p(1-p^{-s})^{-1}$ over the primes converges absolutely,
uniformly on $\operatorname{Re}s\ge1+\delta$ for every $\delta>0$, and equals $\sum_{n\ge1}n^{-s}$
(Euler). More generally, if $\sum_p\abs{b_p}$ converges uniformly on compact sets then
$\prod_p(1+b_p)$ is holomorphic and nonvanishing where no factor vanishes.
\end{enumerate}
\end{theorem}
\noindent\emph{First needed:} \Cref{ch:modforms} (the form of \curve{11a1} is $\eta(z)^2\eta(11z)^2$);
\Cref{ch:lfunctions} (Euler products of $L(E,s)$, $L(\ad f,s)$ and $L(\Sym f,s)$).
\emph{Reference:} \cite[Chapter 5, Section 2]{Ahlfors79}; \cite[Chapter 5]{SteinShakarchi03}.

\begin{theorem}[the Gamma function, complex form (stated; reference)]\label{thm:prereq:gamma}
\leavevmode
\begin{enumerate}[label=(\alph*)]
\item The integral $\Gamma(s)=\int_0^\infty e^{-t}t^{s-1}\,dt$ converges absolutely and locally
uniformly for $\operatorname{Re}s>0$, and $\Gamma$ is holomorphic there.
\item (Functional equation) $\Gamma(s+1)=s\,\Gamma(s)$. Used backwards, it extends $\Gamma$ to a
meromorphic function on $\C$ whose only singularities are simple poles at $s=0,-1,-2,\dots$, with
residue $(-1)^n/n!$ at $s=-n$. $\Gamma$ has no zeros, so $1/\Gamma$ is entire.
\item $\Gamma(s)\Gamma(1-s)=\pi/\sin(\pi s)$; $\Gamma(1/2)=\sqrt\pi$;
$\Gamma(s)\Gamma(s+\tfrac12)=2^{1-2s}\sqrt\pi\,\Gamma(2s)$ (Legendre's duplication formula).
\item (Stirling) In every sector $\abs{\arg s}\le\pi-\delta$,
$\Gamma(s)\sim\sqrt{2\pi}\,s^{s-1/2}e^{-s}$ as $\abs s\to\infty$; on a vertical line,
$\abs{\Gamma(\sigma+it)}\sim\sqrt{2\pi}\,\abs t^{\sigma-1/2}e^{-\pi\abs t/2}$ as $\abs t\to\infty$.
\end{enumerate}
\end{theorem}
\noindent\emph{First needed:} \Cref{ch:lfunctions} (the Gamma factors of $\Lambda(E,s)$ and of the
completed symmetric square $L$-function; the poles of $\Gamma$ account for the trivial zeros).
\emph{Reference:} \cite[Chapter 5, Section 2]{Ahlfors79}; \cite[Chapter 6]{SteinShakarchi03}.

\begin{remark}[the Mellin transform]\label{rem:prereq:mellin}
For $\phi\colon(0,\infty)\to\C$, the \emph{Mellin transform} is
$(\mc M\phi)(s)=\int_0^\infty\phi(t)\,t^{s}\,\frac{dt}{t}$, defined on the vertical strip where the
integral converges absolutely and holomorphic there by \Cref{thm:prereq:improper}(c). By
definition $\Gamma(s)=\mc M(e^{-t})(s)$ for $\operatorname{Re}s>0$, and the substitution
$t\mapsto t/n$ gives $\mc M(e^{-nt})(s)=n^{-s}\Gamma(s)$. Hence, if
$f(iy)=\sum_{n\ge1}a_ne^{-2\pi ny}$ with coefficients of polynomial growth, then for
$\operatorname{Re}s$ large, interchanging sum and integral by absolute convergence,
\[
\int_0^\infty f(iy)\,y^{s}\,\frac{dy}{y}=(2\pi)^{-s}\Gamma(s)\sum_{n\ge1}a_nn^{-s}.
\]
The Mellin transform of a cusp form along the imaginary axis is its $L$-function completed by the
Gamma factor; the transformation law of $f$ under $z\mapsto-1/(Nz)$ becomes, through the
substitution $y\mapsto1/(Ny)$, the functional equation of $\Lambda(f,s)$, and the rapid decay of
$f(iy)$ at both ends gives the analytic continuation. This is the mechanism of \Cref{ch:lfunctions},
and the one by which Riemann continued $\zeta(s)$ using the theta function. For the physics reader:
the Mellin transform is the Fourier transform of \cite[Chapter 13]{RHB06} in the variable $\log t$.
\end{remark}

\section[Group theory and representation theory]{Group theory and representation theory as in Riley, Hobson and Bence}\label{sec:prereq:groups}

This section is what \Cref{ch:groups} and \Cref{ch:reps} assume and then rebuild with proofs.
The reader has met it in \cite[Chapters 28 and 29]{RHB06}, over $\C$ and for finite groups; the
remark at the end says what changes in the book.

\begin{definition}[groups]\label{def:prereq:groups}
A \emph{group} is a set with an associative multiplication, an identity and inverses; it is
\emph{abelian} if the multiplication is commutative, and its \emph{order} is its number of
elements. A \emph{subgroup} is a subset closed under multiplication and inverses; the
\emph{order of an element} $g$ is the order of the cyclic subgroup $\langle g\rangle$ it generates.
A (left) \emph{coset} of $H$ is a set $gH$; the \emph{index} $[G:H]$ is the number of cosets. $H$ is
\emph{normal} if $gHg^{-1}=H$ for all $g$; then the cosets form the \emph{quotient group} $G/H$. A
\emph{homomorphism} $\varphi\colon G\to G'$ satisfies $\varphi(gh)=\varphi(g)\varphi(h)$; its
\emph{kernel} is $\varphi^{-1}(1)$. Two elements are \emph{conjugate} if $h=gxg^{-1}$ for some $g$;
the \emph{conjugacy classes} partition $G$; the \emph{centraliser} $C_G(x)$ is the subgroup of
elements commuting with $x$.
\end{definition}

\begin{theorem}[finite groups (stated; reference)]\label{thm:prereq:groups}
\begin{enumerate}[label=(\alph*)]
\item (Lagrange) The order of a subgroup divides the order of the group; hence the order of every
element divides $\#G$, and $g^{\#G}=1$.
\item The kernel of a homomorphism is a normal subgroup, the image is a subgroup, and
$G/\ker\varphi\iso\im\varphi$.
\item The conjugacy class of $x$ has $\#G/\#C_G(x)$ elements.
\item A group of prime order is cyclic; a cyclic group of order $n$ is isomorphic to $\Z/n\Z$ and
has exactly one subgroup of each order dividing $n$.
\item The examples to keep in mind: $\Z/n\Z$; the symmetric groups $S_n$; the dihedral group $D_n$
of order $2n$ (the symmetries of a regular $n$-gon); the quaternion group
$Q_8=\{\pm1,\pm i,\pm j,\pm k\}$; the matrix groups $\GL_n(\R)$, $\GL_n(\C)$, $\mathrm O(n)$,
$\mathrm{SO}(n)$, $\mathrm U(n)$, $\mathrm{SU}(n)$.
\end{enumerate}
\end{theorem}
\noindent\emph{First needed:} \Cref{ch:groups}, which proves all of this and continues with group
actions, Sylow's theorems, semidirect products, finite abelian groups, profinite groups and the
groups $\GL_2(\F_q)$, which RHB does not treat. \emph{Reference:} \cite[Chapter 28]{RHB06};
\cite[Chapters 1 to 3]{DummitFoote}.

\begin{definition}[representations and characters over $\C$]\label{def:prereq:reps}
A \emph{representation} of $G$ of degree $n$ over $\C$ is a homomorphism $D\colon G\to\GL_n(\C)$
(RHB's ``matrix representation''); equivalently, an action of $G$ by linear maps on an
$n$-dimensional complex vector space $V$, a basis of $V$ turning one into the other. Two
representations are \emph{equivalent} if $D'=P^{-1}DP$ for an invertible $P$ (the same action in
another basis). A subspace $U\subseteq V$ is \emph{invariant} if $gU\subseteq U$ for all $g$; $V$ is
\emph{irreducible} if it has no invariant subspace other than $0$ and $V$, and \emph{completely
reducible} (or \emph{semisimple}) if it is a direct sum of irreducible invariant subspaces. The
\emph{character} of $D$ is $\chi_D(g)=\tr D(g)$: a class function with $\chi_D(1)=n$ and, for $G$
finite, $\chi_D(g^{-1})=\overline{\chi_D(g)}$. Direct sums and tensor products of representations
have characters the sum and the product. The \emph{regular representation} is $V=\C[G]$, with basis
$(e_h)_{h\in G}$ and $g\cdot e_h=e_{gh}$; its character is $\#G$ at $g=1$ and $0$ elsewhere.
\end{definition}

\begin{theorem}[representations of finite groups over $\C$ (stated; reference)]\label{thm:prereq:reps}
Let $G$ be finite.
\begin{enumerate}[label=(\alph*)]
\item (Maschke) Every representation over $\C$, or over any field whose characteristic does not
divide $\#G$, is completely reducible: every invariant subspace has an invariant complement.
\item (Schur) A $G$-equivariant linear map between irreducible representations is $0$ or an
isomorphism, and the $G$-equivariant endomorphisms of an irreducible complex representation are the
scalars.
\item (Orthogonality) For irreducible characters $\chi,\chi'$,
$\frac1{\#G}\sum_{g\in G}\chi(g)\overline{\chi'(g)}$ is $1$ if $\chi=\chi'$ and $0$ otherwise. The
irreducible characters form an orthonormal basis of the space of class functions, so their number
is the number of conjugacy classes; and $\sum_\chi\chi(1)^2=\#G$.
\item A complex representation is determined up to equivalence by its character; the multiplicity
of the irreducible $\chi$ in a representation of character $\psi$ is
$\frac1{\#G}\sum_g\psi(g)\overline{\chi(g)}$; a representation is irreducible if and only if
$\frac1{\#G}\sum_g\abs{\chi(g)}^2=1$. The regular representation contains each irreducible
representation with multiplicity equal to its degree.
\item (Column orthogonality) $\sum_\chi\chi(g)\overline{\chi(h)}$ is $\#C_G(g)$ if $g$ and $h$
are conjugate and $0$ otherwise.
\item If $G$ is abelian, every irreducible complex representation has degree $1$; the degree-one
representations (``characters'' in the second sense) of $G$ form a group isomorphic to $G$, and for
$G=\Z/n\Z$ they are $k\mapsto e^{2\pi ijk/n}$, $0\le j<n$. The degree of an irreducible
representation of any finite $G$ divides $\#G$.
\end{enumerate}
\end{theorem}
\noindent\emph{First needed:} \Cref{ch:reps}. \emph{Reference:} \cite[Chapter 29]{RHB06};
\cite[Part I]{SerreLinearReps}; \cite[Chapter 18]{DummitFoote}.

\begin{remark}[what changes in the book]\label{rem:prereq:changes}
Four things change when the paper's representations arrive, all in \Cref{ch:reps}.
(1) \emph{The field.} The representations of the paper live over $\F_\ell$, $\Zl$, $\Ql$, their
algebraic closures, and rings of integers of number fields. When the characteristic $\ell$ divides
the order of the group, as it does for $\GL_2(\F_3)$ in characteristic $3$ and for an inertia group
acting on $E[\ell]$ through a unipotent element of order $\ell$, Maschke's theorem fails: an
invariant subspace need not have an invariant complement, and a representation need not be
semisimple. The invariants are then the composition factors (Jordan-Holder), not direct summands,
and a character determines only the semisimplification. The right statement is Brauer-Nesbitt: two
semisimple representations with the same characteristic polynomial of $\rho(g)$ for every $g$ are
isomorphic; the traces determine the characteristic polynomials in characteristic $0$ or when the
characteristic exceeds the dimension, but not in general.
(2) \emph{Irreducible versus absolutely irreducible.} Over a field that is not algebraically closed
an irreducible representation can become reducible after extending scalars, and Schur's lemma in
the form ``equivariant endomorphisms are scalars'' holds only for absolutely irreducible ones. The
paper's example: the cyclic group of order $4$ generated by $\sm01{-1}0$ acts irreducibly on
$\F_3^2$ ($X^2+1$ has no root in $\F_3$) but reducibly on $\F_9^2$, as $\psi\oplus\psi^{-1}$ for a
character $\psi$ of order $4$. This is why the hypothesis of the paper's main theorem says
``absolutely irreducible''.
(3) \emph{The group.} The Galois groups $\GQ$, $\GQp$, $I_p$ are infinite profinite groups, their
representations are required to be continuous, and the finite theory is applied through finite
quotients.
(4) \emph{Lattices.} A representation on a $\Ql$-vector space $V$ with a stable $\Zl$-lattice $T$
can be reduced modulo $\ell$ to a representation on $T/\ell T$ over $\F_\ell$; the reduction depends
on $T$, its semisimplification does not, and when the reduction is irreducible all stable lattices
are homothetic (Ribet's lemma).
Finally, the word ``character'' in the book almost always means a degree-one representation of a
group, usually a profinite one (the cyclotomic character $\chil$, the quadratic characters
$\chi_{-3}$, $\chi_{-4}$, $\chi_{q^*}$, Dirichlet characters), not a trace function; \Cref{ch:reps} says
``the character of $\rho$'' when the trace is meant.
\end{remark}

\section{Elementary probability and statistics}\label{sec:prereq:stats}

This section is used only in \Cref{ch:computation}, which explains the paper's data protocol
(Section sec:data of the paper) and its calibration argument (Section sec:sharp). The reader's
course in statistics is more than enough; the cards fix the vocabulary.

\begin{definition}[finite probability spaces; base rates]\label{def:prereq:prob}
A \emph{finite probability space} is a finite set $\Omega$ with a function $P\colon\Omega\to[0,1]$
of total sum $1$; an \emph{event} is a subset $A\subseteq\Omega$, with $P(A)=\sum_{\omega\in A}P(\omega)$;
a \emph{random variable} is a function $X\colon\Omega\to\R$, with \emph{expectation}
$E[X]=\sum_\omega X(\omega)P(\omega)$ and \emph{variance} $\operatorname{Var}X=E[(X-E[X])^2]$. The
\emph{indicator} $1_A$ of an event has $E[1_A]=P(A)$. Events $A,B$ are \emph{independent} if
$P(A\cap B)=P(A)P(B)$, and the \emph{conditional probability} is $P(A\mid B)=P(A\cap B)/P(B)$. When
$\Omega$ is a population with the uniform distribution $P(\omega)=1/\#\Omega$, $P(A)=\#A/\#\Omega$
is the \emph{proportion} of $A$ in the population, also called the \emph{base rate} of the
attribute $A$, and $E[X]$ is the population mean of $X$. If $A$ and $B$ were independent, the
proportion of $A$ within $B$ would equal the base rate of $A$; comparing the two is the simplest
test of association.
\end{definition}

\begin{theorem}[expectation (stated; reference)]\label{thm:prereq:expectation}
\begin{enumerate}[label=(\alph*)]
\item (Linearity) $E[aX+bY]=aE[X]+bE[Y]$, with no assumption of independence. In particular the
expected number of events among $A_1,\dots,A_n$ that occur is $\sum_iP(A_i)$, whether or not the
$A_i$ are independent.
\item For independent $X,Y$: $E[XY]=E[X]E[Y]$ and $\operatorname{Var}(X+Y)=\operatorname{Var}X+\operatorname{Var}Y$.
\item (Markov, Chebyshev) For $X\ge0$ and $t>0$, $P(X\ge t)\le E[X]/t$; for any $X$,
$P(\abs{X-E[X]}\ge t)\le\operatorname{Var}X/t^2$.
\item (Law of large numbers) If $X_1,\dots,X_n$ are independent with the same distribution, of mean
$m$ and finite variance, then $P(\abs{(X_1+\dots+X_n)/n-m}\ge\eps)\to0$ as $n\to\infty$, for every
$\eps>0$.
\end{enumerate}
\end{theorem}
\noindent\emph{First needed:} \Cref{ch:computation}. \emph{Reference:} \cite[Chapters I, V and
IX]{Feller68}; \cite[Chapter 30]{RHB06}.

\begin{definition}[hypothesis tests]\label{def:prereq:test}
A \emph{hypothesis test} consists of a \emph{null hypothesis} $H_0$, a precise probabilistic model
for the data (for instance: two attributes are independent); a \emph{test statistic} $T$, a number
computed from the data; and a \emph{rejection rule}, fixed before the data are examined, of the form
``reject $H_0$ if $T\in R$'', where $R$ is chosen so that $P_{H_0}(T\in R)\le\alpha$ for a
prescribed \emph{significance level} $\alpha$. The \emph{$p$-value} of an observed value $t$ is
$P_{H_0}(T\text{ at least as extreme as }t)$. Rejecting $H_0$ says that the data would be very
unlikely under $H_0$; it does not say that a particular alternative is true, and it says nothing at
all if the rule was chosen after looking at the data.
\end{definition}

\begin{example}[the paper's calibration argument]\label{ex:prereq:calibration}
The population is the set of working-set curves ($N\le300000$) with $E[3]$ irreducible and without
complex multiplication. Among the $366151$ of them for which the right-hand side of the paper's
rule (its inequality eq:rule) is $0$, the base rate of $v_3(\degphi)=0$ is $53.2\%$. Take as null
hypothesis that $v_3(\degphi)$ is independent of the local data. Then a curve whose bound is $b$
would have $v_3(\degphi)<b$ with the probability that a bound-$0$ curve has $v_3(\degphi)<b$, and
summing these probabilities over the $867578$ working-set curves with a positive bound gives, by
\Cref{thm:prereq:expectation}(a), an expected number of about $605000$ violations. The observed
number is $0$. The force of the argument lies in the comparison of these two numbers, not in a
$p$-value (which would be astronomically small): the inequality is not an artefact of most modular
degrees being divisible by $3$. What the argument does not show is that the inequality holds beyond
the data; that is the business of the theorem under its hypotheses and of the conjecture outside
them. The equality rates ($57.5\%$ at $\ell=3$) and the mean excess ($0.63$) that the paper reports
are descriptive statistics, proportions and means in the population, not tests.
\end{example}

\begin{definition}[hold-out sets; pre-registration]\label{def:prereq:holdout}
When a data set is searched for a relation among many candidates, some candidate will fit the data
by chance, and the more thorough the search the more surely (the problem of \emph{multiple
comparisons}); a relation found on a data set therefore cannot be tested on the same data. The
remedy is to split the data, before the search, into a \emph{working set}, used for the search and
for everything else, and a \emph{hold-out set}, which is not examined; to write down the statement
and the test, with its rejection rule, before looking at the hold-out (\emph{pre-registration}); and
then to test once. A statement that passes such a test has been \emph{blind-tested}. A statement
formulated after the hold-out has been used can only be \emph{confirmed} on it, which is weaker,
and a fresh hold-out is needed to blind-test it.
\end{definition}

\begin{paperlink}\label{pl:prereq:protocol}
The paper's ranges are the working set ($N\le300000$), the hold-out ($300000<N\le400000$), the
extension ($400000<N\le500000$) and the out-of-table curves ($500000<N\le4\cdot10^6$). The
component-group inequality (eq:main) and Conjecture conj:eis were found on the working set, written
down, tested once on the hold-out, then on the extension, and nothing was adjusted afterwards; the
twist-symmetrised form (eq:twist) was likewise tested once on the hold-out and the extension. The
Euler-factor rule (Conjecture conj:euler) was formulated last, after the hold-out and the extension
had been used, so for it they are a confirmation and the out-of-table curves, generated before the
rule was formulated, are the blind test. \Cref{ch:computation} explains each range and each count
of the paper's tables. \emph{Reference:} \cite{Wasserman04}; \cite[Chapter 31]{RHB06}.
\end{paperlink}

\section{Where to look}\label{sec:prereq:books}

\emph{Real analysis.} Rudin \cite{Rudin76} is the reference used above and contains everything in
\Cref{sec:prereq:real} with proofs; Apostol \cite{Apostol74} treats double series, improper
integrals and the Gamma function in more detail. RHB \cite[Chapter 4]{RHB06} covers series and
convergence tests at the level of a first course.

\emph{Complex analysis.} Ahlfors \cite{Ahlfors79} and Stein-Shakarchi \cite{SteinShakarchi03} both
contain all of \Cref{sec:prereq:complex}; the second has a chapter on the Gamma and zeta functions
and one on the upper half plane. The first chapter of Diamond-Shurman \cite{DiamondShurman} does the
upper half plane, the action of $\SL_2(\Z)$ and $q$-expansions in exactly the form
\Cref{ch:modforms} needs. RHB \cite[Chapters 24 and 25]{RHB06} covers residues and contour
integration.

\emph{Linear algebra.} Axler \cite{Axler15} (over $\R$ and $\C$, with the spectral theorem and
generalised eigenvectors done carefully) and Halmos \cite{Halmos74} (bilinear forms, duality,
tensor products, in a coordinate-free way) between them contain \Cref{sec:prereq:linear} over
$\R$ and $\C$. For arbitrary fields, and for the same constructions over rings, use Lang
\cite[Chapters XIII to XVI]{Lang} or Dummit-Foote \cite[Chapters 11 and 12, and Section
10.4]{DummitFoote}. RHB \cite[Chapter 8]{RHB06} is the reader's starting point.

\emph{Groups and representations.} RHB \cite[Chapters 28 and 29]{RHB06} is what we assume. Serre
\cite{SerreLinearReps} is the standard reference: Part I for characteristic $0$, Part III for the
modular theory that \Cref{ch:reps} needs. Dummit-Foote \cite[Chapters 18 and 19]{DummitFoote} and
Lang \cite[Chapter XVIII]{Lang} cover the same ground, and \cite[Chapters 1 to 6]{DummitFoote} is
the book's own reference for the group theory of \Cref{ch:groups}.

\emph{Probability and statistics.} RHB \cite[Chapters 30 and 31]{RHB06} is what we assume. Feller
\cite{Feller68} is the classical text on discrete probability; Wasserman \cite{Wasserman04} is a
concise account of statistical inference, including hypothesis tests and multiple comparisons, at
the level needed to read \Cref{ch:computation}.

\emph{After this appendix.} The book's own references for the material it builds are named in the
``Notes and references'' section of each chapter; the ones used most are Lang \cite{Lang} and
Dummit-Foote \cite{DummitFoote} for algebra, Atiyah-Macdonald \cite{AtiyahMacdonald} for
commutative algebra, Serre \cite{SerreLocalFields} and Neukirch \cite{Neukirch} for local fields and
number fields, Silverman \cite{SilvermanAEC,SilvermanATAEC} for elliptic curves, and
Diamond-Shurman \cite{DiamondShurman} for modular forms.
